\documentclass[11pt,a4paper]{article}
\usepackage{fontspec,amsmath,amssymb,amsthm,microtype}
\setmainfont{lmroman10-regular.otf}[BoldFont=lmroman10-bold.otf,ItalicFont=lmroman10-italic.otf,BoldItalicFont=lmroman10-bolditalic.otf]
\usepackage[margin=25mm]{geometry}
\usepackage[unicode,hidelinks]{hyperref}
\hypersetup{pdfauthor={LangevinLean Project}}
\setlength{\emergencystretch}{3em}
\setlength{\parskip}{0.4em}
\setlength{\parindent}{0pt}
\allowdisplaybreaks[2]
\raggedbottom
\title{\texorpdfstring{From summed Gaussian residuals to the exact Dynkin identity}{From summed Gaussian residuals to the exact Dynkin identity}}
\author{LangevinLean Project}
\date{}
\begin{document}\maketitle
Draft mathematical exposition; these notes may contain mistakes.


This is the deterministic limit argument connecting the proved Gaussian weak steps to a continuous integral solution. It does not assume a stochastic calculus identity or a Wasserstein error recurrence.



For \(h_n=T/(n+1)\), the residual hypothesis is arbitrarily small at first order:


\[\forall\varepsilon>0\ \exists K\ge0:\quad
|r_n|\le C\bigl(\varepsilon+K\sqrt{h_n}\bigr)
\ \text{eventually},\qquad r_n\longrightarrow0.\]


Matching-noise Euler convergence and uniform continuity of the two observations then give the exact identity


\[\operatorname{test}(Y(T))-\operatorname{test}(Y(0))
=\int_0^T\operatorname{generator}(Y(t))\,dt.\]
\section{\texorpdfstring{Force the accumulated weak residual to zero}{Force the accumulated weak residual to zero}}


Suppose every positive \(\varepsilon\) permits an eventual bound \(C[\varepsilon +K\sqrt{T/(n+1)}]\), with \(K\) allowed to depend on \(\varepsilon\). First fix \(\varepsilon\) small relative to the desired error; then refine the mesh so its square-root term is equally small. This proves that the actual scalar residual converges to zero.


\section{\texorpdfstring{Identify the unique limiting residual}{Identify the unique limiting residual}}


Uniform continuity sends the terminal matching-noise Euler state to the terminal test observation. Uniform mesh convergence sends its generator sum to the time integral. Their difference therefore converges to the proposed Dynkin residual, while the summed weak remainder converges to zero. Uniqueness of real limits gives the exact identity.



Earlier results used here: \cite[Section~3.16.1]{LangevinLean2026ULASC}; \cite[Section~3.16.8]{LangevinLean2026ULASC}.

\begin{thebibliography}{ULA-LSI}
\bibitem[ULA-SC]{LangevinLean2026ULASC}
LangevinLean Project. \emph{Unadjusted Langevin under Strong Convexity}. LL-01, 2026. Draft report, content revision 263fb89cdbf0. \url{https://odakyildiz.github.io/LangevinLean/docs/pdf/UnadjustedLangevin/FullProof.pdf}.
\end{thebibliography}

\end{document}
